%-------------------------------------------------------------------------------
% §4 PROJETS DE GRANDE ENVERGURE PERTINENTS
%-------------------------------------------------------------------------------
\cvsection{4. Projets de grande envergure pertinents}

\cvgroupreset

\begin{cvgroup}
\cvsubsection{Projet 1 : OBR / MFE / ASYCUDA / EBMS | Nemidis Technologies}

\begin{cvmetadata}
\textbf{Période :} Novembre 2025 – Juin 2026 \\
\textbf{Rôle :} Architecte logiciel / Expert en intégration de systèmes — supervision technique \\
\textbf{Client :} Office Burundais des Recettes (OBR) — Institution publique, système d'information national
\end{cvmetadata}

\begin{cvprojectsummary}
Supervision et mise en œuvre de l'intégration entre ASYCUDA, EBMS et POS-OBR pour assurer la traçabilité des marchandises importées depuis le dédouanement jusqu'à la facturation électronique au sein du système d'information de l'OBR. Modernisation du SI fiscal national et digitalisation des processus publics dans le cadre de la stratégie d'e-gouvernement du Burundi.
\end{cvprojectsummary}

\begin{cvdetailslist}{Réalisations}
  \item {Définition de l'architecture d'intégration et rédaction des spécifications fonctionnelles et techniques (cahier des charges ASYCUDA–EBMS–POS-OBR).}
  \item {Modélisation UML par diagrammes de séquences des flux inter-systèmes et des interactions entre acteurs métier.}
  \item {Conception et supervision du module de gestion des marchandises importées : stock douane, validation contribuable, synchronisation POS local.}
  \item {Mise en œuvre de la traçabilité basée sur les références douanières DMC et intégration opérationnelle ASYCUDA–EBMS–POS-OBR.}
  \item {Coordination fonctionnelle et technique entre Nemidis, les équipes OBR et les parties prenantes institutionnelles.}
  \item {Arbitrage des choix technologiques, validation des livrables et préparation du déploiement sur infrastructure cloud.}
  \item {Estimation des charges, planification du projet et suivi de la conformité aux exigences réglementaires et non-fonctionnelles (sécurité, traçabilité).}
\end{cvdetailslist}

\cvdetailblock{Technologies}{UML (diagrammes de séquences) · Dart · Flutter · Docker · Google Cloud Platform · PostgreSQL}

\cvdetailblock{Compétences démontrées}{Institution publique · SI national · E-gouvernement · Architecture cible · Cahier des charges · Normes de développement · Intégration de systèmes · Supervision · Coordination institutionnelle · Gouvernance des données publiques}
\end{cvgroup}

\begin{cvgroup}
\cvsubsection{Projet 2 : Plateforme nationale d'ordonnances électroniques | Doctolib}

\begin{cvmetadata}
\textbf{Période :} Juin 2020 – Décembre 2022 \\
\textbf{Rôle :} Ingénieur logiciel senior — conception et intégration \\
\textbf{Client :} Téléservice public de santé numérique — intégration aux services numériques de l'État français
\end{cvmetadata}

\begin{cvprojectsummary}
Conception, développement et maintenance de l'application d'ordonnances électroniques, téléservice public de santé numérique utilisé quotidiennement par les professionnels de santé en France et intégré au système national d'e-Prescription.
\end{cvprojectsummary}

\begin{cvdetailslist}{Réalisations}
  \item {Conception et développement full-stack du module d'ordonnances médicales (Ruby on Rails, React) à forte volumétrie sur base PostgreSQL.}
  \item {Pilotage de l'intégration du système national d'e-Prescription : sécurisation et vérification des ordonnances via QR Code.}
  \item {Intégration avec les API et services numériques de l'État français dans le cadre de la dématérialisation des parcours de soins.}
  \item {Garantie de l'intégrité, de la sécurité et de la traçabilité des données médicales sur un SI national critique.}
  \item {Validation de la conformité des livrables avec les cahiers des charges et les normes fonctionnelles et non-fonctionnelles (sécurité, accessibilité, réglementation).}
  \item {Participation aux déploiements continus (Docker) et à la maintenance d'un système d'information national.}
  \item {Formation et mentorat de nouveaux ingénieurs sur les bonnes pratiques de qualité logicielle.}
\end{cvdetailslist}

\cvdetailblock{Technologies}{Docker · Ruby on Rails · React · PostgreSQL · API REST (JSON)}

\cvdetailblock{Compétences démontrées}{Secteur public numérique · SI national · Interopérabilité · Intégration API · Sécurité des données · Qualité et conformité · Services publics numériques · Renforcement des capacités}
\end{cvgroup}

\begin{cvgroup}
\cvsubsection{Projet 3 : Plateforme d'intégration de données financières | Climb}

\begin{cvmetadata}
\textbf{Période :} Janvier 2023 – Juin 2025 \\
\textbf{Rôle :} Ingénieur logiciel senior — architecture et intégration \\
\textbf{Client :} Plateforme numérique de conseil patrimonial et fiscal
\end{cvmetadata}

\begin{cvprojectsummary}
Conception et évolution d'une plateforme d'intégration de données financières reliant cinq compagnies d'assurance (Vie Plus, Generali, Apicil, Swiss Life, Lombard) au système d'information central. Collecte, transformation, stockage et synchronisation des données pour alimenter les applications métier, Salesforce et les espaces clients.
\end{cvprojectsummary}

\begin{cvdetailslist}{Réalisations}
  \item {Architecture d'une plateforme d'intégration multi-partenaires : collecte automatisée des données de cinq assureurs.}
  \item {Mise en place de flux d'échanges via API REST (JSON), SFTP, webhooks et fichiers XML.}
  \item {Normalisation des données financières dans un modèle unique PostgreSQL et gouvernance de la qualité des données.}
  \item {Développement d'API REST (JSON) consommées par les applications web et les espaces clients.}
  \item {Synchronisation des contrats, mouvements financiers, allocations et valorisations entre assureurs, Salesforce et applications internes.}
  \item {Supervision des flux de données, contrôles de qualité et amélioration continue des performances.}
  \item {Participation aux choix d'architecture, aux déploiements et à l'évolution de la plateforme.}
\end{cvdetailslist}

\cvdetailblock{Technologies}{PostgreSQL · Redis · Ruby on Rails · API REST (JSON)}

\cvdetailblock{Compétences démontrées}{Intégration de SI · Architecture de données · API REST (JSON) · Gouvernance des données · Synchronisation multi-systèmes · Plateformes digitales · Supervision technique}
\end{cvgroup}

\begin{cvgroup}
\cvsubsection{Projet 4 : Lingu.Africa — plateforme multilingue d'apprentissage | Initiative personnelle}

\begin{cvmetadata}
\textbf{Période :} Mars 2021 – Présent \\
\textbf{Rôle :} Fondateur / Architecte applicatif \\
\textbf{Contexte :} Inclusion numérique et éducation — diaspora africaine
\end{cvmetadata}

\begin{cvprojectsummary}
Plateforme numérique créée pour transmettre le \textbf{Kirundi} aux enfants de la diaspora burundaise, aujourd'hui étendue à \textbf{26 langues africaines}. Conception, développement et exploitation d'une solution scalable de diffusion de contenus éducatifs multilingues.
\end{cvprojectsummary}

\begin{cvdetailslist}{Réalisations}
  \item {Définition de l'architecture globale, du modèle de données et des choix technologiques de la plateforme.}
  \item {Développement full-stack (Next.js, React) : commerce électronique, gestion de contenu et workflows de publication automatisés.}
  \item {Montée en charge progressive de 1 langue (Kirundi) à 26 langues africaines.}
  \item {Pilotage du déploiement, de la supervision et de l'évolution continue de la solution.}
\end{cvdetailslist}

\cvdetailblock{Technologies}{Next.js · React · PostgreSQL}

\cvdetailblock{Compétences démontrées}{Plateforme digitale scalable · Inclusion numérique · Kirundi · Ownership bout-en-bout · Architecture applicative · Développement web moderne}
\end{cvgroup}
